% ---------------------------------------------------------------------------
% Font macros for the F Ring User's Guide
%
% The Word original uses several distinctive fonts that are not generally
% available to LaTeX.  Each is wrapped in a macro here so the font choice can
% be changed in ONE place.  Current substitutions:
%
%   Word font           Used for                          LaTeX substitute
%   ------------------  --------------------------------  -----------------------
%   Times New Roman     body text, title                  Times (mathptmx)
%   Calibri (bold)      headings                          Helvetica (phv)
%   Cascadia Code       file names, program constants,    Courier (pcr)
%                       code blocks, directory listings
%   Courier New         "Literal" style (rarely used)     Courier (pcr)
%   Lucida Sans         (label listings in older drafts;  Helvetica (phv)
%                       retained for completeness)
%
% If you compile on a full TeX system you may prefer, e.g.:
%   \renewcommand{\codefontfamily}{\fontfamily{zi4}\selectfont}   % Inconsolata
% or with LuaLaTeX/XeLaTeX + fontspec, set these to the real fonts.
% ---------------------------------------------------------------------------

% --- inline font for file names / program constants (Word: Cascadia Code)
\newcommand{\codefontfamily}{\fontfamily{pcr}\selectfont}
\newcommand{\code}[1]{{\codefontfamily #1}}

% --- inline font for the Word "Literal" character style (Word: Courier New)
\newcommand{\literalfontfamily}{\fontfamily{pcr}\selectfont}
\newcommand{\literal}[1]{{\literalfontfamily #1}}

% --- inline font for XML label excerpts in older drafts (Word: Lucida Sans)
\newcommand{\labelxmlfontfamily}{\fontfamily{phv}\selectfont}
\newcommand{\labelxml}[1]{{\labelxmlfontfamily #1}}

% --- heading font (Word: Calibri bold)
\newcommand{\headingfont}{\sffamily\bfseries}

% --- title font (Word: Times New Roman 18 pt bold)
\newcommand{\usgtitlefont}{\rmfamily\bfseries\fontsize{18pt}{22pt}\selectfont}

% --- Zotero citation placeholder.
% The Word document manages citations with Zotero.  Until the bibliography is
% converted to BibTeX (second pass), \zcite simply prints the plain citation
% text captured from the Word field.  The full CSL JSON data for every
% citation was extracted to citations/zotero-citations.json.
\newcommand{\zcite}[1]{#1}

% --- verbatim environments for code blocks --------------------------------
% CodeBlock    : Word "Code" style paragraphs (Cascadia Code)
% LiteralBlock : Word "Literal" style paragraphs (Courier New)
% LabelBlock   : XML label listings (Lucida Sans in older drafts)
% A [fontsize=...] option is emitted by the converter where the Word source
% used a smaller point size (e.g. the long XML label listings).
% breaklines wraps lines too long for the page, as Word does; the
% continuation line is indented and unmarked (set breaksymbolleft to change).
\DefineVerbatimEnvironment{CodeBlock}{Verbatim}%
  {fontsize=\small,formatcom=\codefontfamily,xleftmargin=1.5em,
   breaklines=true,breakanywhere=true,breaksymbolleft={},breakindent=2em}
\DefineVerbatimEnvironment{LiteralBlock}{Verbatim}%
  {fontsize=\small,formatcom=\literalfontfamily,xleftmargin=1.5em,
   breaklines=true,breakanywhere=true,breaksymbolleft={},breakindent=2em}
\DefineVerbatimEnvironment{LabelBlock}{Verbatim}%
  {fontsize=\small,formatcom=\labelxmlfontfamily,xleftmargin=1.5em,
   breaklines=true,breakanywhere=true,breaksymbolleft={},breakindent=2em}
